% 可供整体重写时纳入的证据章节；未接入旧 main.tex，旧 PDF 不代表当前实现。
% 数字与实验条件见 docs/results/2026-10-28-efficiency/README.md。
\chapter{受限文件系统的恢复、查询与软件写入开销}
\label{chap:current-evidence}

本章以 xv6 的教学内核为对象，研究固定挂载点上的文件操作、版本化对象图及可重复的软件路径计数。它不是完整 POSIX 文件系统，也不是对物理掉电的一般性证明。文件内容经 KV run 保存，并保留内容定义分块（CDC）与 stream；这些组合及其变体不应被表述为对已有 COW、持久内存文件系统或过滤索引的首创。持久内存文件系统和字节寻址一致性的相关研究可参见 \cite{pmfs2014,nova2016,bpfs2009}。

\section{实现边界与设计}

VFS 保留原生 xv6 根，将 \texttt{/pm} 作为固定挂载点，按组件解析路径；目录引用与普通文件描述符固定底层 inode。独立打开同一 PM inode 时共用工作缓冲而各自维护偏移；\texttt{dup}/\texttt{fork} 的同一打开对象共享偏移。当前没有动态挂载、PM 硬链接、打开后 unlink、符号链接、跨后端重命名、权限或 \texttt{mmap}；文件大小上限 1024 字节，运行路径仍由全局 PMFS 锁串行化。

元数据使用不可变 inode 和 dentry 快照，双槽 checkpoint 指向完整对象图。非替换重命名只改变 dentry 图，因而追加一个新 dentry 页并复用已提交的 inode 页；替换已有目标仍需同时产生两页。文件内容同步先发布新 KV run，再发布新 inode 快照，并复用未改变的 dentry 根。该设计要求旧根始终指向不可变且未回收的页；本原型未实现安全垃圾回收，分配空间最终耗尽。

普通文件先预探测 CDC 切分：当块数不超过四个 extent 且不多于 256 字节固定分块时采用 CDC，否则以固定分块兜底。此策略保证有界文件可表示，并不实现内容去重，也不保证所有输入的 CPU 开销降低。KV 冷查询首先查有界精确索引；超过 512 个已跟踪键后，未跟踪键使用完整的 newest-first fallback。coarse posting 对所有已提交键维护，fallback 倒序枚举置位的 run，再经过 Bloom 和 fence 判断。更新已经跟踪的键即使在索引满后也必须刷新指向。

\section{可复现的软件请求消融}

下列数据来自 2026 年 10 月 28 日的单核 QEMU 9.2.0 TCG、EMUL 一次指定测试工作负载；补丁仅分别关闭对应优化，测试程序同步修改预期计数。\texttt{dev\_write\_bytes} 是软件接口复制与 checkpoint 状态字 store 的计数，不是实际介质写入；run probe 是访问候选 run 描述符的次数，不是磁盘读取次数。数值和逐次日志均需以仓库记录复核，不能外推为通用吞吐或介质寿命。

\begin{table}[htbp]
\centering
\caption{相同指定工作负载中的软件请求计数}
\label{tab:current-software-efficiency}
\begin{tabular}{lrr}
\toprule
路径与单位 & 对照 & 当前实现 \\
\midrule
非替换重命名，4 KiB 页写请求 & 2 & 1 \\
同上，软件写入字节 & 8256 & 4160 \\
4 字节文件内容同步，4 KiB 页写请求 & 3 & 2 \\
同上，软件写入字节 & 12416 & 8320 \\
旧 run 中的冷键，run probe & 144 & 0 \\
旧 run 中的冷键，查询页读取 & 1 & 1 \\
稀疏 posting 旧键，run probe & 80 & 16 \\
未使用的 posting 桶，run probe / 页读取 & 214 / 0 & 0 / 0 \\
\bottomrule
\end{tabular}
\end{table}

在指定测试中，相同的 4 字节覆盖写跳过了设备写请求；已存在的脏状态没有被清除。1024 字节固定模式与四块 CDC 模式都经内容读回验证。容量测试中，完整元数据对照成功重命名 1014 次，本实现成功 2030 次；这是一次有界、追加分配工作负载的结果，不等于耐久寿命提高。

\section{恢复实验与尚未解决的问题}

三核共享内存后端分别执行非替换与替换子树重命名的 43 次启动，以及文件/KV 的 62 次启动；实验性 QEMU NVDIMM 后端及其 Zicbom 指令路径分别执行同样的测试。检查器覆盖指定的半页、整页与 checkpoint 软件切点、旧/新内容及多次重启。模拟未覆盖任意指令位置或物理缓存失效；普通共享文件与杀死 QEMU 进程也不能证明宿主断电后的持久性。有界规格模型枚举旧/新结果并验证负控，但尚未从真实内核提取完整 store/flush 轨迹。

剩余的软件问题包括安全回收和 compaction、多个仍受保护的 checkpoint 的可达性、固定 checkpoint 槽的热点、可扩展并发以及更完整的 VFS 语义。软件页写、flush 行请求不能称为 NAND 擦除、器件磨损或真实介质写放大；没有真实设备证据时不得声称已实现磨损均衡或提升了硬件寿命。原生 xv6 对照也没有与 PMFS 等价的设备持久化边界，故现有 TCG tick 不能支持硬件性能优越性结论。
